\chapter{Artificial Intelligence in Contemporary Biomedical Science}

\section{Introduction}

Artificial intelligence (AI) has rapidly transformed into a central methodological framework within modern biomedical science. In the last decade, AI applications in biomedicine have expanded to encompass stem cell differentiation modeling, target identification for gene therapy, microbial image synthesis, vaccine epitope design, protein structure prediction, and clinical decision support systems, among many others \parencite{topol2019, jumper2021}. This expansion has been driven by three converging developments: exponential growth of multimodal biomedical datasets, maturation of deep learning architectures, and widespread availability of high-performance computational resources.

Recent work in cellular therapy demonstrates how AI can model dynamic gene regulatory networks to guide stem cell differentiation, integrating synthetic biology with predictive modeling \parencite{choudhury2025}. Similarly, computational biology and machine learning approaches have enhanced workflows in gene therapy, including target identification, vector design, and outcome prediction \parencite{danaeifar2025}. In structural biology, the development of AlphaFold2 \parencite{jumper2021} marked a crucial moment, demonstrating that deep learning could solve one of the field's oldest problems --- the protein folding problem --- with near-experimental accuracy. More recently, generative approaches have been applied to design novel proteins and nucleic acid constructs with defined functional properties \parencite{ferruz2022}.

Despite these advances, the rapid expansion of AI methodologies has raised important theoretical and practical questions that the scientific community has not yet fully resolved. Central among these are the following: Are current AI systems genuinely explanatory, or are they sophisticated pattern recognizers without mechanistic insights? Can deep learning models be trusted in high-risk clinical contexts where errors carry life-altering consequences? How should uncertainty be quantified, propagated, and communicated to end users? And to what extent do biases encoded in training data compromise the generalizability and fairness of AI-guided recommendations?

These questions collectively define the contemporary methodological debate in biomedical AI. The present chapter examines this debate through the major paradigms --- supervised and unsupervised learning, deep learning and representation learning, generative models and synthetic data, Bayesian and probabilistic inference, large language models and hybrid systems and AutoML --- with particular attention to both demonstrated scientific potential and unresolved challenges. Our treatment emphasizes epistemic responsibility: the responsibility of researchers not only to understand what their models predict, but why, under what conditions, and with what degree of confidence.

\section{Supervised and Unsupervised Machine Learning in Biomedicine}

\subsection{Supervised Learning: Applications and Scope}

Methods of supervised learning --- including support vector machines (SVMs), random forests, gradient boosted trees, and feedforward neural networks --- have become central to analyses in genomics, transcriptomics, proteomics, and clinical informatics \parencite{libbrecht2015}. In genomics, supervised classifiers trained on labeled sequence data have been applied to identification of pathogenic single nucleotide polymorphisms (SNPs), prediction of splice sites, and classification of regulatory elements. In transcriptomics, machine learning models trained on RNA-seq data enable differential expression analysis, cell type deconvolution, and prediction of phenotype at the transcriptome level.

In gene therapy research, supervised learning has been applied to all phases of the therapeutic pipeline. At the target identification stage, models trained on curated disease-gene association databases --- such as DisGeNET or the OMIM catalog --- predict novel therapeutic targets from multi-omic profiles \parencite{danaeifar2025}. At the vector design stage, classifiers predict the efficiency and specificity of adenoassociated virus (AAV) serotypes for particular tissue tropism based on capsid sequence features. At the outcome prediction stage, models trained on patient clinical and genomic data estimate therapeutic response, risk of adverse event, and long-term durability of gene correction.

Proteomics has similarly benefited from supervised learning approaches. Mass spectrometry data processed through machine learning pipelines now enable high-throughput protein identification, quantification, and mapping of post-translational modifications, enabling large-scale profiling of disease-associated proteomes \parencite{orre2019}. In clinical imaging, supervised convolutional neural networks (CNNs) achieved performance comparable to radiologists on tasks such as tumor detection, classification, and segmentation across multiple cancer types \parencite{litjens2017}.

\subsection{Unsupervised Learning and Discovery of Latent Structure}

Unsupervised learning techniques --- including clustering algorithms, dimensionality reduction methods, and hierarchical modeling --- enable discovery of latent biological structure in high-dimensional omic datasets without the need for pre-labeled training data. Principal component analysis (PCA), t-distributed stochastic neighbor embedding (t-SNE), and Uniform Manifold Approximation and Projection (UMAP) have become standard tools for visualization of single-cell RNA-seq (scRNA-seq) data, enabling researchers to identify transcriptionally distinct cell populations with unprecedented resolution \parencite{mcinnes2018, becht2018}.

Weighted gene co-expression network analysis (WGCNA) exemplifies a biologically motivated unsupervised approach that reconstructs networks of co-regulated genes from expression data, associating network modules with clinical phenotypes or experimental conditions \parencite{langfelder2008}. These module-trait associations have been used to nominate gene regulatory hubs as therapeutic targets across a range of diseases, including neurodegenerative disorders, cancer, and metabolic disease.

In the context of cellular therapy, unsupervised methods have been applied to characterize the heterogeneity of stem cell populations, identify rare subpopulations with distinct differentiation potentials, and map trajectories of cell state transitions during reprogramming \parencite{saelens2019}. Integration of multimodal single-cell data --- combining transcriptomics, chromatin accessibility (scATAC-seq), and proteomics --- presents ongoing challenges for unsupervised analysis, motivating the development of new methods for dimensionality reduction and data integration such as MOFA+ and Seurat v4 \parencite{argelaguet2020, hao2021}.

\subsection{Debates: Generalizability, Bias, and Uncertainty}

The debate surrounding supervised and unsupervised approaches centers on three interrelated concerns: generalizability across populations and contexts, encoding of systematic bias, and inadequate quantification of predictive uncertainty. Many biomedical machine learning models are trained on datasets derived from non-representative population samples --- often predominantly European ancestry --- which substantially limits their external validity and clinical utility in underrepresented groups \parencite{popejoy2016, obermeyer2019}. This concern is particularly acute for polygenic risk scores, which have been demonstrated to exhibit substantial performance disparities across ancestry groups when derived from eurocentric genome-wide association study (GWAS) cohorts \parencite{martin2019}.

Performance metrics in published studies of biomedical machine learning frequently emphasize aggregate accuracy, area under the receiver operating characteristic curve (AUROC), or related summary statistics, without sufficient attention to calibration --- the degree to which predicted probabilities correspond to empirical event frequencies. Well-calibrated uncertainty estimates are essential for clinical deployment, where a model expressing 90\% confidence in an incorrect diagnosis carries very different implications from one expressing 55\% confidence \parencite{guo2017}. Prospective validation in diverse and independent cohorts under realistic operational conditions remains the exception rather than the rule in published literature \parencite{wynants2020}.

\section{Deep Learning and Representation Learning}

\subsection{Architectures and Biomedical Applications}

Deep learning has enabled automated and hierarchical extraction of features from complex, high-dimensional biological data, fundamentally altering research workflows in imaging, genomics, and structural biology. Convolutional neural networks (CNNs) have been applied to automated histopathological image analysis, demonstrating that deep architectures can identify morphological features predictive of cancer subtype, prognosis, and treatment response at a scale that would be impractical for human annotators \parencite{esteva2017, kather2019}. In hematology, CNNs trained on microscopy images of blood cells enable automated classification of hematopoietic cell types and identification of abnormal populations relevant to leukemia diagnosis \parencite{matek2019}.

In the domain of cellular therapy, deep learning supports quality control of stem cell cultures by detecting morphological deviations from expected differentiation trajectories in time-lapse image data. Mesenchymal stem cell (MSC) preparations, critical for a range of regenerative medicine applications, have been characterized using deep CNNs trained to predict immunomodulatory potency from phase-contrast microscopy images, reducing dependence on costly and time-consuming functional assays \parencite{kusumoto2021}.

Recurrent architectures, including long short-term memory (LSTM) networks and transformer-based models, have been applied to genomic sequence modeling, prediction of gene expression from chromatin state, and analysis of clinical time series. The transformer architecture, originally developed for natural language processing, has been adapted for genomics in models such as Enformer \parencite{avsec2021}, which predicts gene expression and chromatin accessibility from DNA sequence with substantially improved long-range interaction modeling compared to previous convolutional architectures.

A notable application of representation learning in microbiology involves generation of synthetic microbial colony datasets using deep learning style transfer \parencite{pawlowski2022}. By transferring visual style features from real colony images to synthetically generated backgrounds, the authors were able to substantially expand training corpora for automated colony detection models without requiring additional manual annotation, while maintaining competitive detection performance. This work illustrates how representation learning techniques can directly address the widespread bottleneck of limited labeled data in biomedical image analysis.

\subsection{Interpretability and the Black Box Problem}

Despite these impressive empirical achievements, deep neural networks remain largely opaque to mechanistic interpretation. The representations learned and encoded in hidden layers of a deep network are not, in general, directly mappable to scientifically meaningful constructs, making it difficult to understand why a given prediction was made or to identify the biological features driving model behavior. This interpretability deficit is a central concern in clinical translation, where regulatory agencies, clinicians, and patients may reasonably demand explanations for algorithmic recommendations \parencite{rudin2019}.

A range of post-hoc interpretability methods have been developed to partially address this limitation, including gradient-based saliency maps, layer-wise relevance propagation (LRP), integrated gradients, and SHAP values (SHapley Additive exPlanations) \parencite{lundberg2017, sundararajan2017}. These methods provide approximate attributions of model predictions to input features, but are not without limitations. Saliency-based methods have been shown to be sensitive to model architecture and initialization, producing inconsistent attributions across ostensibly equivalent models \parencite{adebayo2018}. Furthermore, attribution of features at the input level does not necessarily illuminate the internal computational mechanisms by which predictions are generated.

There is growing recognition within the field that intrinsically interpretable models --- which are designed to be transparent by construction, rather than explained post hoc --- may be preferable to black box alternatives in high-risk domains, even at some cost to predictive performance \parencite{rudin2019}. The debate continues regarding the degree to which interpretability and predictive accuracy are fundamentally in tension, and whether advances in architecture design can produce models that are both powerful and mechanistically transparent.

\section{Generative Models and Synthetic Data}

\subsection{Generative Adversarial Networks and Variational Autoencoders}

Generative models introduced a new paradigm in biomedical AI, enabling synthesis of realistic data samples, design of novel molecular entities, and expansion of training datasets. Generative adversarial networks (GANs), introduced by \textcite{goodfellow2014}, have been applied to synthesis of histopathological images, expansion of MRI datasets, generation of realistic electronic health record (EHR) data, and de novo design of molecular structures with desired physicochemical properties \parencite{kadurin2017, baowaly2019}.

Variational autoencoders (VAEs) provide an alternative generative framework that imposes a structured probabilistic prior on the latent space, enabling controlled generation and interpretable interpolation between data points. VAEs have been applied in drug discovery to navigate chemical space toward regions predicted to exhibit desired properties of potency, selectivity, and pharmacokinetics \parencite{gomez2018}. Hybrid architectures combining the representational power of GANs with the latent structure of VAEs have further expanded the generative toolkit available to biomedical researchers.

\subsection{Protein and Nucleic Acid Design}

Perhaps the most scientifically significant application of generative models in biomedicine is in structural and functional protein design. Building on the representational foundations established by AlphaFold2 \parencite{jumper2021}, generative protein language models such as ProtGPT2 \parencite{ferruz2022} and ESMFold \parencite{lin2023} have enabled design of novel protein sequences with predicted structural properties. These tools reduce the combinatorial search space of protein engineering from an astronomically large sequence space to a tractable design problem, with implications for enzyme engineering, therapeutic antibody design, and synthetic biology.

In vaccine development, generative and predictive modeling are increasingly integrated across the vaccine design pipeline. AI-guided epitope prediction tools identify immunogenic peptide candidates from pathogen proteomes, while mRNA structure optimization algorithms design sequences that maximize translational efficiency and stability for mRNA vaccine platforms \parencite{imani2025}. Diffusion model-based approaches --- analogous to those used in image generation --- have been applied to design protein linkers with specified target interfaces, representing a cutting-edge application of generative AI in therapeutic protein engineering \parencite{watson2023}.

\subsection{Synthetic Data: Promises and Risks}

Synthetic data generation addresses a critical bottleneck in biomedical AI: the scarcity of large, annotated, and privacy-compliant datasets. Privacy regulations --- including the Health Insurance Portability and Accountability Act (HIPAA) in the United States and the General Data Protection Regulation (GDPR) in the European Union --- impose substantial restrictions on sharing and secondary use of patient data, limiting the scale and diversity of training corpora available to researchers. Synthetic data generated by GAN or diffusion models can, in principle, be shared without direct privacy risk while preserving the statistical properties of the underlying clinical population.

Nevertheless, synthetic datasets are not a panacea. Generative models trained on non-representative source data will faithfully reproduce, and may even amplify, the biases present in those data. Evaluation of synthetic data quality requires careful comparison to real data distributions in clinically meaningful dimensions, and there is evidence that naively evaluated realism may mask systematic distortions in rates of rare events, minority subgroup characteristics, and outcome distributions \parencite{walonoski2018}. Furthermore, the use of synthetic data for regulatory submission or clinical validation remains an evolving area, with guidance from agencies such as the FDA still being formalized \parencite{fda2023}.

\section{Bayesian Networks and AI-Guided Probabilistic Inference}

\subsection{Principled Modeling of Uncertainty}

Bayesian approaches offer a mathematically principled framework for modeling uncertainty in biomedical AI, integrating prior domain knowledge with observed data through Bayes' theorem to produce posterior probability distributions over parameters and predictions. Unlike point-estimate discriminative models, Bayesian inference explicitly represents and propagates uncertainty, making it well-suited to high-risk biomedical applications where the cost of overconfident predictions is substantial \parencite{ghahramani2015}.

In gene therapy, Bayesian networks have been applied to model the probabilistic dependencies among genomic, transcriptomic, and clinical variables, enabling principled target prioritization under uncertainty \parencite{danaeifar2025}. In epidemiology and pharmacovigilance, hierarchical Bayesian models are used to estimate adverse event rates from spontaneous reporting systems, sharing statistical strength across related drug-event combinations while appropriately representing the uncertainty associated with rare outcomes. In vaccine development, Bayesian adaptive trial designs use interim data to update prior beliefs about vaccine efficacy and safety, enabling more efficient use of clinical trial resources \parencite{berry2010}.

\subsection{Structure Learning and Causal Inference}

Beyond its utility for uncertainty quantification, Bayesian networks offer a framework for representation and learning of causal dependencies among variables --- a property that aligns more closely with the explanatory goals of scientific investigation than purely associative models. Causal models based on directed acyclic graphs (DAGs) encode hypothesized causal relationships among biological variables and can, under appropriate assumptions about data-generating mechanisms, be used to predict the consequences of experimental interventions \parencite{pearl2009}.

Causal discovery algorithms --- such as the PC algorithm, FCI, and more recent variants augmented with deep learning --- attempt to infer causal graph structure from observational data, a notoriously challenging problem requiring strong assumptions \parencite{spirtes2000}. In single-cell genomics, causal discovery has been applied to inference of gene regulatory networks from scRNA-seq data, with the goal of identifying target-transcription factor relationships that can be experimentally validated and therapeutically exploited \parencite{skokgibbs2022}.

\subsection{Hybrid Probabilistic-Deep Learning Models}

The contemporary debate in this area centers on whether hybrid models that combine deep learning representations with probabilistic graphical model inference can achieve superior interpretability and uncertainty quantification without sacrificing the representational power that has made deep learning so empirically successful. Approaches such as deep kernel learning, neural process models, and Bayesian neural networks represent attempts to integrate these paradigms \parencite{wilson2020}. Gaussian process regression, a Bayesian nonparametric approach, has been applied to model gene expression dynamics as a function of pseudotime in single-cell trajectory analysis, providing principled uncertainty bounds on inferred differentiation trajectories \parencite{lonnberg2017}. The practical challenge of scaling Bayesian inference to the parameter regimes of modern deep learning architectures --- often involving billions of parameters --- remains an active area of research.

\section{Large Language Models in Scientific Workflows}

\subsection{Applications in Biomedical Research}

Large language models (LLMs) --- including GPT-4 \parencite{openai2023}, Gemini \parencite{googledeepming2023}, and domain-adapted models such as BioMedLM, BioGPT \parencite{luo2022}, and Med-PaLM 2 \parencite{singhal2023} --- are increasingly integrated into scientific workflows. These models, trained on vast corpora of biomedical literature, clinical notes, and genomic sequences, exhibit emergent capabilities including literature synthesis, hypothesis generation, answering clinical questions, and automation of bioinformatics pipelines. In biomedical named entity recognition (NER), relation extraction, and clinical note processing, LLMs have achieved state-of-the-art performance on established benchmarks, often without task-specific fine-tuning \parencite{wei2022}.

Protein language models (PLMs) represent a related family of models trained not on natural language but on large databases of amino acid sequences. Models such as ESM-2 \parencite{lin2023} and ProtTrans \parencite{elnaggar2022} learn contextualized representations of amino acid positions that encode evolutionary, structural, and functional information, and have been used to predict the effects of mutations on protein stability and function, classify protein families, and guide the generation of novel sequences. The success of PLMs illustrates how the transformer architecture originally developed for language can be productively adapted to biological sequence data.

\subsection{Epistemic Limitations and Risks}

Despite their capabilities, LLMs exhibit significant limitations that are particularly consequential in biomedical contexts. Hallucination --- the generation of plausible but factually incorrect or fabricated information --- is a well-documented failure mode of current autoregressive language models \parencite{ji2023}. In the context of scientific literature synthesis, LLMs have been observed generating plausible-appearing but nonexistent citations, attributing fabricated claims to legitimate-appearing authors and journals. This poses a non-trivial risk to research integrity when LLM outputs are incorporated into manuscripts or systematic reviews without independent verification.

LLMs also lack explicit epistemic grounding: they do not maintain formal representations of knowledge or inference mechanisms that distinguish known facts from plausible interpolations of training data. They cannot reliably distinguish between well-established scientific consensus and contested or preliminary findings, and may reproduce outdated information or misconceptions widely discussed in the training corpus. Clinical deployment of LLMs in decision support roles therefore requires robust safeguards, including retrieval-augmented generation (RAG) architectures that ground model outputs in verified and up-to-date knowledge bases, uncertainty-aware output calibration, and mandatory human-in-the-loop review \parencite{lewis2020}.

Regulatory and ethical frameworks for LLM deployment in clinical settings are still nascent. The FDA has issued design guidance on AI-enabled devices but has not yet comprehensively addressed LLM-specific risks \parencite{fda2023}. Questions of accountability when an LLM contributes to a clinical error, informed consent for AI-assisted decision-making, and equitable access to LLM-powered tools across healthcare settings remain largely unresolved.

\section{Hybrid Systems and AutoML}

\subsection{Hybrid AI: Combining Paradigms}

Hybrid AI systems integrate complementary paradigms --- symbolic reasoning, probabilistic graphical models, and deep learning --- with the goal of retaining the strengths of each while mitigating their individual weaknesses. Neurosymbolic AI approaches, for example, seek to combine the pattern recognition capabilities of deep networks with the rule-based compositional reasoning of symbolic AI, thereby enabling models that can learn from data while also respecting logical constraints and providing interpretable reasoning traces \parencite{marcus2019, garcez2023}.

In biomedical applications, hybrid systems have been developed for clinical decision support, integrating deep learning models for medical image analysis with knowledge graph-based reasoning over clinical ontologies such as SNOMED CT and the Gene Ontology. These architectures can ground predictions in structured biological knowledge, reducing the risk of predictions that are statistically plausible but biologically incoherent. Hybrid approaches are also prominent in drug discovery, where graph neural networks (GNNs) operating on molecular graphs are combined with physics-based force fields and quantum chemical calculations to predict binding affinities and ADMET properties with improved reliability \parencite{wieder2020}.

\subsection{AutoML: Democratization and Its Pitfalls}

Automated machine learning (AutoML) platforms --- including AutoKeras, TPOT, H2O AutoML, and cloud-hosted offerings from major technology vendors --- lower barriers to model development by automating the processes of feature engineering, algorithm selection, hyperparameter optimization, and architecture search. In biomedical research contexts where deep expertise in machine learning may not be available, AutoML enables researchers to deploy competitive predictive models without manual tuning, potentially accelerating discovery and broadening access to AI-assisted analysis \parencite{he2021}.

However, the democratization provided by AutoML raises serious concerns regarding methodological transparency and scientific reproducibility. Researchers who deploy AutoML pipelines without fully understanding the underlying algorithmic choices --- the selection of a particular learner, the application of specific preprocessing transformations, or the configuration of cross-validation procedures --- may inadvertently violate key modeling assumptions or introduce subtle forms of data leakage. Published studies based on AutoML analyses may be difficult to reproduce if the AutoML system introduces non-determinism or if the specific pipeline configuration is not completely reported. There is risk that AutoML adds even greater speed to the existing tendency to report predictive performance metrics without adequate mechanistic interpretation or external validation \parencite{kapoor2023}.

\subsection{Federated Learning and Privacy-Preserving AI}

An important variant of hybrid and distributed AI is federated learning, which enables model training across multiple institutions without centralizing sensitive patient data \parencite{rieke2020}. In federated learning, each participating institution trains a local model on its own data, and only model parameters --- rather than raw data --- are communicated to a central aggregation server. This architecture can substantially expand the effective training set available to biomedical AI models without violating patient privacy regulations and is particularly valuable for rare diseases where no single institution has sufficient cases for robust model training. However, federated learning introduces its own technical challenges, including heterogeneity of data distributions across sites (the so-called ``non-IID'' data), communication overhead, and susceptibility to adversarial model poisoning attacks \parencite{kairouz2021}.

\section{Persistent Challenges}

The preceding sections have identified a constellation of challenges that cut across AI paradigms and biomedical application domains. Taken collectively, they define the frontier of contemporary debate.

\textbf{Interpretability and mechanistic understanding.} The opacity of deep learning models remains the most frequently cited barrier to clinical translation. While post-hoc interpretability methods provide partial mitigation, they do not fully resolve the underlying problem: that many high-performing models are not designed to be mechanistically informative. Advances in intrinsically interpretable architectures, as well as theoretical understanding of what neural networks actually compute, are needed \parencite{rudin2019}.

\textbf{Reproducibility and methodological rigor.} A substantial proportion of published biomedical AI studies have been shown to suffer from inadequate reporting of methodological details, inappropriate data-splitting practices, overfitting to benchmark datasets, and failure to validate on independent prospective cohorts \parencite{kapoor2023, wynants2020}. The development of field-specific reporting standards --- analogous to CONSORT for clinical trials and STROBE for observational studies --- has been proposed as a partial remedy \parencite{mongan2020}.

\textbf{Bias, fairness, and equity.} Systematic underrepresentation of diverse populations in biomedical training datasets has direct consequences for the fairness and generalizability of AI-derived clinical tools. Addressing this requires not only technical interventions --- such as reweighting, adversarial debiasing, and equity-constrained optimization --- but also structural changes to data collection practices and research infrastructure \parencite{obermeyer2019}.

\textbf{Regulatory oversight and clinical validation.} The regulatory landscape for AI-enabled devices is rapidly evolving, but significant gaps remain. Continuous learning systems --- which update their parameters in response to new data after initial deployment --- present particular regulatory challenges, as post-market model drift may render an approved device progressively less safe or effective without triggering re-evaluation \parencite{fda2023}. International harmonization of regulatory standards is also needed as medical AI devices are increasingly deployed across jurisdictions.

\textbf{Epistemological tension.} Perhaps more fundamentally, there is an unresolved epistemological tension between the goal of predictive performance and the goal of scientific understanding. A model that accurately predicts patient outcomes without any mechanistic interpretability may be clinically useful but scientifically unsatisfying --- and may fail silently in situations that differ from its training distribution in ways that a mechanistic model could detect. Biomedical sciences have historically prioritized causal explanation over predictive empiricism, and there is legitimate debate regarding whether the broad adoption of black box AI represents a departure from that epistemic tradition that should be embraced or resisted \parencite{lipton2018}.

\section{Conclusion}

Artificial intelligence has catalyzed transformative advances across gene therapy, cellular therapy, vaccine design, structural biology, and computational biomedicine more broadly. The empirical achievements documented across supervised learning, deep learning, generative models, Bayesian inference, and language models collectively demonstrate that AI is not merely a computational convenience but a genuine engine of scientific discovery, capable of identifying patterns, generating hypotheses, and accelerating experimental workflows at a scale and speed that transcends unaided human cognitive capacity.

However, predictive performance alone is insufficient for sustainable scientific and clinical integration. The biomedical community faces a constellation of unresolved challenges --- interpretability, reproducibility, bias, regulatory adequacy, and epistemological coherence --- that collectively demand a more deliberate and critically rigorous approach to AI methodology adoption. Deploying powerful but opaque models without adequate validation infrastructure, reporting standards, or governance frameworks risks a form of methodological excess that could ultimately undermine confidence in AI-assisted science.

The path forward requires coordinated effort across multiple dimensions. Technical research must advance interpretable architectures, principled uncertainty quantification, and robust methods for detecting distributional shift. Clinical validation must evolve to prioritize prospective, multi-site evaluation in diverse patient populations over retrospective benchmarking in convenience cohorts. Reporting standards must mandate transparent disclosure of methodological details sufficient for independent replication. Regulatory frameworks must adapt to the distinct challenges posed by continuous learning systems and LLM-based decision support. And the scientific community must engage in sustained reflection on the epistemological role of AI in biomedical investigation --- not as a substitute for mechanistic understanding, but as a powerful complement to it.

The contemporary debate is therefore not whether AI should be used in biomedical science, but how it should be integrated into scientific methodology in a manner that preserves --- and ultimately deepens --- our epistemic standards for what it means to understand biological systems and provide safe, effective, and equitable clinical care.

\clearpage
\begin{revise}\small
\begin{enumerate}[itemsep=0pt, topsep=0pt, parsep=0pt]
    \item Supervised learning: SVMs, random forests, gradient boosting, feedforward networks; applications to pathogenic SNPs, splice site prediction, regulatory element classification, AAV vector design, proteomics (mass spectrometry).
    \item Unsupervised learning: PCA, t-SNE, UMAP for scRNA-seq; WGCNA for co-expression networks; MOFA+, Seurat v4 for multimodal integration.
    \item Debates: generalizability (European ancestry biases in GWAS/polygenic scores), calibration vs.\ AUROC, rare prospective validation \parencite{popejoy2016,obermeyer2019}.
    \item Deep learning: CNNs (histopathology, hematopoietic cell classification), LSTMs, Transformers; Enformer (gene expression from DNA sequence) \parencite{avsec2021}.
    \item Interpretability: ``black box''; post-hoc methods (saliency maps, LRP, SHAP, integrated gradients) are unstable \parencite{adebayo2018}; intrinsically interpretable models preferable in high-risk contexts \parencite{rudin2019}.
    \item Generative models: GANs (Goodfellow, 2014) for histopathological images and synthetic EHR; VAEs for chemical space navigation in drug discovery; diffusion models for protein linker design.
    \item Protein design: AlphaFold2 \parencite{jumper2021}, ProtGPT2 \parencite{ferruz2022}, ESMFold; vaccine epitope prediction; mRNA optimization.
    \item Synthetic data: address scarcity + privacy (HIPAA/GDPR), but amplify source biases; FDA regulation evolving.
    \item Bayesian networks: principled uncertainty; DAGs for causal inference (PC algorithm, FCI); hybrid probabilistic--deep models (deep kernel, neural processes, Bayesian neural networks).
    \item LLMs: GPT-4, BioGPT, Med-PaLM 2; hallucination = central risk; lack epistemic grounding; mitigation via RAG + human-in-the-loop.
    \item AutoML (AutoKeras, TPOT, H2O): democratizes ML but double black box; risk of data leakage and irreproducibility. Federated learning: trains without centralizing data; challenges = non-IID data, adversarial poisoning.
\end{enumerate}
\end{revise}

\clearpage

\printbibliography[heading=subbibliography,title={References}]

